%% Experimental setup: datasets, the disjoint observation-process splits, metrics, protocol.
%% Included by sigconf.tex. Do not add a preamble here.

\section{Experimental Setup}
\label{sec:setup}

\subsection{Datasets and tasks}
We use five datasets spanning three sensing modalities, two tasks and two backbones
(Table~\ref{tab:datasets}). Preprocessing, splits and normalisation are the unmodified published pipelines:
t-PatchGNN \cite{zhang2024tpatchgnn} for the forecasting benchmarks (PhysioNet from the 2012 challenge
\cite{silva2012physionet}) and Raindrop \cite{zhang2022raindrop} for classification, including the
minority-class upsampling of the binary tasks. USHCN appears at two horizons because horizon length changes how
much a model must rely on its history, which turns out to change its sensitivity to the shift.

\begin{table}[t]
\caption{Datasets and task settings. Backbone follows the published protocol per task.}
\label{tab:datasets}
\small
\begin{tabular}{llrrll}
\toprule
setting & task & vars & samples & backbone & metric \\
\midrule
PhysioNet 24h$\to$24h & forecast & 41 & 12{,}000 & iTransformer & MSE \\
Activity 3000$\to$1000\,ms & forecast & 12 & 5{,}400 & iTransformer & MSE \\
USHCN 24$\to$1 mo & forecast & 5 & 26{,}736 & iTransformer & MSE \\
USHCN 24$\to$12 mo & forecast & 5 & 26{,}736 & iTransformer & MSE \\
P12 mortality & classify (2) & 36 & 11{,}988 & PatchTST & AUROC \\
P19 sepsis & classify (2) & 34 & 38{,}803 & PatchTST & AUROC \\
PAM activity & classify (8) & 17 & 5{,}333 & PatchTST & accuracy \\
\bottomrule
\end{tabular}
\end{table}

\subsection{Observation-process splits}
\label{sec:splits}
Training and evaluation draw from \emph{disjoint} observation-process distributions, so that generalisation across
observation processes is measured rather than assumed.

\textbf{Training pool} $P_{\mathrm{mix}}$ is deliberately mild: random dropout at $p\in\{0.1,0.25,0.4\}$, stride
$2$, burst at $p\in\{0.2,0.3\}$, variable dropout at $p=0.2$ and local masking at $p=0.3$.

\textbf{Evaluation conditions} are the clean condition plus 15 shifted processes, each applied by the
re-observation protocol of Section~\ref{sec:reobs}. Nine are designated \emph{unseen} because their family never
appears in the training pool: random dropout $0.7$, variable dropout $0.5$, stride $3$ and $4$, long gap $0.4$,
burst $0.5$, variable-group dropout $0.3$, and two compositions (random $0.5$ $+$ variable $0.3$; stride $3$ $+$
variable-group $0.3$). The remaining six re-use training families at severities outside the pool, together with
density-controlled sampling at $\rho=0.5$. We report degradation over all fifteen shifted conditions and,
separately, over the nine unseen ones.

\textbf{Mining candidate pool.} Two pool constructions are reported. $\mathcal{P}=P_{\mathrm{mix}}$ is the natural
default: mining over the processes the model already trains on. The \emph{harsher} pool adds severities and
pairwise compositions beyond $P_{\mathrm{mix}}$ while remaining disjoint from the evaluation conditions, so no
evaluation process is ever mined.

\subsection{Metrics}
Forecasting uses masked MSE and MAE over observed target entries, averaged per variable and then over variables
with supervision. Binary classification uses AUROC and AUPRC; the eight-class task uses accuracy and macro F1. We
summarise sensitivity by the relative degradation of a model with respect to its \emph{own} clean score,
\begin{equation}
\Delta_{\mathrm{shift}}=\frac{1}{|\mathcal{E}|}\sum_{\Sobs\in\mathcal{E}}
\frac{\mathrm{err}_{\Sobs}-\mathrm{err}_{\mathrm{clean}}}{\mathrm{err}_{\mathrm{clean}}},
\label{eq:degradation}
\end{equation}
over the shifted conditions $\mathcal{E}$, and report the same quantity restricted to the unseen subset. For
classification we take $\mathrm{err}=1-\mathrm{AUROC}$ (or $1-\mathrm{accuracy}$) so that $\Delta$ is computed on
the same metric used to judge the task; cross-entropy degradation can move in the opposite direction to AUROC and
we do not use it as a headline number.

Two cautions govern how we read Eq.~\ref{eq:degradation}. It is normalised by a model's own clean score, so a
uniformly weaker model can post a flattering ratio, and it says nothing about attainable error. We therefore report
clean, mean shifted and worst-case absolute performance in every table and treat $\Delta_{\mathrm{shift}}$ as a
secondary summary that must be read beside them. We also report parameter count, training time per epoch, peak
training memory and inference latency.

\subsection{Protocol}
All models use the query-based irregular embedding with hidden width 64, four heads and three layers, trained with
Adam at $10^{-3}$, batch 32 (128 for USHCN; 16 for P12, where a single view's attention already occupies
4.7\,GB), an epoch cap of 40, early-stopping patience 10 and a minimum of 15 epochs. Configurations that are
compared are trained at identical batch size. Runs execute sequentially on one 8\,GB GPU, and every run records
its configuration, git commit, seed, checkpoint and metrics. \textbf{All reported numbers use a single seed} and
are marked \prelim; the protocol calls for five paired seeds, and we regard nothing reported here as a statistical
result. \todo{five paired seeds for the primary comparisons}
